
	\documentclass[spanish, 11pt]{exam}
	
	%These tell TeX which packages to use.
	\usepackage{array,epsfig}
	\usepackage{amsmath, textcomp}
	\usepackage{amsfonts}
	\usepackage{amssymb}
	\usepackage{amsxtra}
	\usepackage{amsthm}
	\usepackage{mathrsfs}
	\usepackage{color}
	\usepackage{multicol}
	\usepackage{verbatim}
	
	
	\usepackage[utf8]{inputenc}
	\usepackage[spanish]{babel}
	\usepackage{eurosym}
	
	\usepackage{graphicx}
	\graphicspath{{../img/}}
	
	
	
	\printanswers
	\nopointsinmargin
	\pointformat{}
	
	%Pagination stuff.
	%\setlength{\topmargin}{-.3 in}
	%\setlength{\oddsidemargin}{0in}
	%\setlength{\evensidemargin}{0in}
	%\setlength{\textheight}{9.in}
	%\setlength{\textwidth}{6.5in}
	%\pagestyle{empty}
	
	\renewcommand{\solutiontitle}{\noindent\textbf{Sol:}\enspace}
	
	\newcommand{\samedir}{\mathbin{\!/\mkern-5mu/\!}}
	
	\newcommand{\class}{1º Bachillerato}
	\newcommand{\examdate}{\today}
	
	\newcommand{\tipo}{A}
	
	
	\newcommand{\timelimit}{50 minutos}
	
	
	
	\pagestyle{head}
	\firstpageheader{\includegraphics[width=0.2\columnwidth]{header_left}}{\textbf{Departamento de Matemáticas\linebreak \class}\linebreak \examnum}{\includegraphics[width=0.1\columnwidth]{header_right}}
	\runningheader{\class}{\examnum}{Página \thepage\ of \numpages}
	\runningheadrule
	
	\newcommand{\examnum}{4 - Polinomios}
    \begin{document}
    \begin{questions}
    \question p012e03 - Dados los polinomios $ A(x)=2{x^3} - 5{x^2} + 6, \  B(x)=- \frac{1}{2}{x^5} - {x^4} + 6x$ halla:
        \begin{multicols}{3} 
        \begin{parts} \part[]  $ A(x) + B(x) $  \begin{solution}  $ - \frac{x^{5}}{2} - x^{4} + 2 x^{3} - 5 x^{2} + 6 x + 6 $  \end{solution} \part[]  $ A(x) - B(x) $  \begin{solution}  $ \frac{x^{5}}{2} + x^{4} + 2 x^{3} - 5 x^{2} - 6 x + 6 $  \end{solution}
        \end{parts}
        \end{multicols}
        \question p012e04 - Dados los polinomios $ A(x)=3{x^3} - 6{x^2} + 2x - 1, \  B(x)=- {x^4} + {x^3} + x - 6, \  C(x)={x^4} - {x^2} + x + \frac{1}{2}$ halla:
        \begin{multicols}{3} 
        \begin{parts} \part[]  $ A(x) \cdot B(x) $  \begin{solution}  $ - 3 x^{7} + 9 x^{6} - 8 x^{5} + 6 x^{4} - 25 x^{3} + 38 x^{2} - 13 x + 6 $  \end{solution} \part[]  $ A(x) - 3B(x) + 5C(x) $  \begin{solution}  $ 8 x^{4} - 11 x^{2} + 4 x + \frac{39}{2} $  \end{solution} \part[]  $ x^2\cdot A(x) + 3x\cdot B(x) $  \begin{solution}  $ - 3 x^{4} + 2 x^{3} + 2 x^{2} - 18 x $  \end{solution}
        \end{parts}
        \end{multicols}
        \question p012e07 - Halla el cociente y el resto de:
        \begin{multicols}{3} 
        \begin{parts} \part[]  $ ( {5{x^4} - 7{x^2} + 6x + 1} ):( {3{x^2}} ) $  \begin{solution}  $ \left ( \frac{5 x^{2}}{3} - \frac{7}{3}, \quad 6 x + 1\right ) $  \end{solution} \part[]  $ ( {7{x^4} - 3{x^2} + 6x - 1} ):( {{x^2}-x+3} ) $  \begin{solution}  $ \left ( 7 x^{2} + 7 x - 17, \quad - 32 x + 50\right ) $  \end{solution} \part[]  $ ( {{x^6} - 5} ):( {{x^2} - x}) $  \begin{solution}  $ \left ( x^{4} + x^{3} + x^{2} + x + 1, \quad x - 5\right ) $  \end{solution} \part[]  $ ( {8{x^6} - 5x^4 + 6 } ):( {2x^2} - 1) $  \begin{solution}  $ \left ( 4 x^{4} - \frac{x^{2}}{2} - \frac{1}{4}, \quad \frac{23}{4}\right ) $  \end{solution} \part[]  $ ( {3{x^5} - 6{x^2} + 9} ):( {{x^2} + 1} ) $  \begin{solution}  $ \left ( 3 x^{3} - 3 x - 6, \quad 3 x + 15\right ) $  \end{solution} \part[]  $ ( {x^9} - 7{x} + 1 ):( {{x^3} + x} ) $  \begin{solution}  $ \left ( x^{6} - x^{4} + x^{2} - 1, \quad - 6 x + 1\right ) $  \end{solution}
        \end{parts}
        \end{multicols}
        \question p012e08 - Dados $ A(x)= - {x^3} + 2{x^2} + 5, \  B(x)= 2{x^4} + 3x + 6$ halla el valor numérico de ambos polinomios en:
        \begin{multicols}{3} 
        \begin{parts} \part[]  $ x=1 $  \begin{solution}  $ 6\ y \ 11 $  \end{solution} \part[]  $ x=-1 $  \begin{solution}  $ 8\ y \ 5 $  \end{solution} \part[]  $ x=2 $  \begin{solution}  $ 5\ y \ 44 $  \end{solution} \part[]  $ x=-2 $  \begin{solution}  $ 21\ y \ 32 $  \end{solution} \part[]  $ x=\frac{1}{2} $  \begin{solution}  $ \frac{43}{8}\ y \ \frac{61}{8} $  \end{solution} \part[]  $ x=-\frac{1}{2} $  \begin{solution}  $ \frac{45}{8}\ y \ \frac{37}{8} $  \end{solution}
        \end{parts}
        \end{multicols}
        \question p012e09 - Halla, para cada uno  de los siguientes polinomios, sus raíces:
        \begin{multicols}{3} 
        \begin{parts} \part[]  $ x^2 - 1 $  \begin{solution}  $ \left\{-1, 1\right\} $  \end{solution} \part[]  $ x^2 - 7 $  \begin{solution}  $ \left\{- \sqrt{7}, \sqrt{7}\right\} $  \end{solution} \part[]  $ 3x^2 -12 $  \begin{solution}  $ \left\{-2, 2\right\} $  \end{solution} \part[]  $ 5x^2 -25 $  \begin{solution}  $ \left\{- \sqrt{5}, \sqrt{5}\right\} $  \end{solution}
        \end{parts}
        \end{multicols}
        \question p012e10 - ¿Tiene el polinomio A(x)= $x^4 +3$ alguna raíz real?
        \begin{multicols}{3} 
        \begin{parts} \part[]  $ x^4 +3 $  \begin{solution}  $ \left\{- \frac{\sqrt{2} \sqrt[4]{3}}{2} - \frac{\sqrt{2} \sqrt[4]{3} i}{2}, - \frac{\sqrt{2} \sqrt[4]{3}}{2} + \frac{\sqrt{2} \sqrt[4]{3} i}{2}, \frac{\sqrt{2} \sqrt[4]{3}}{2} - \frac{\sqrt{2} \sqrt[4]{3} i}{2}, \frac{\sqrt{2} \sqrt[4]{3}}{2} + \frac{\sqrt{2} \sqrt[4]{3} i}{2}\right\} $  \end{solution}
        \end{parts}
        \end{multicols}
        
    \end{questions}
    \end{document}
    